\begin{abstract}
O monitoramento microbiológico de superfícies na indústria depende hoje de plaqueamento, que leva
dias para gerar resultado. Este trabalho apresenta uma plataforma que combina biossensores
eletroquímicos impressos (eletrodos de carbono com sonda de captura para rRNA 16S), aquisição por
voltametria de pulso diferencial (DPV) e aprendizado de máquina para classificar superfícies como
contaminadas ou não contaminadas e estimar a carga bacteriana. Os voltamogramas são lidos
diretamente dos arquivos do potenciostato, passam por controle de qualidade automático e são
organizados em um banco de dados rastreável. A característica mais informativa foi a curva de
diferença entre estágios do sensor no mesmo eletrodo (antes e depois do contato com o alvo).
Foram comparadas \NCombinacoes{} combinações de pré-processamento e modelo por validação cruzada
agrupada por amostra; o modelo escolhido (\MelhorModelo, \MelhorVariante) obteve acurácia
balanceada de \ABcv{} na validação cruzada e de \ABteste{} no conjunto de teste
(\TipoTeste), com AUC de \AUCteste. A carga bacteriana das amostras contaminadas foi estimada com
RMSE de \RegRMSEteste{} log$_{10}$\,UFC/mL. \preencher{revisar após o experimento com os dados
reais das empresas validadoras}

\medskip\noindent\textbf{Palavras-chave:} biossensor eletroquímico; voltametria de pulso
diferencial; rRNA 16S; aprendizado de máquina; monitoramento de superfícies; ISO 18593.
\end{abstract}
